% DIAPOSITIVA 5: Pregunta central
\begin{frame}{Pregunta de investigación}
  \begin{block}{Pregunta central}
    ¿Es posible diseñar un modelo de \textit{nowcasting} electoral \textbf{elección-agnóstico},
    cuyos parámetros se entrenen con múltiples contiendas pasadas y sean
    \textbf{transferibles} a elecciones nuevas, que opere \textbf{sin encuestas},
    usando exclusivamente señales de redes sociales?
  \end{block}

  \vspace{0.25em}
  \begin{columns}[T]
    \begin{column}{0.48\textwidth}
      {\small\textbf{\color{barBlue}Preguntas derivadas (P)}}
      \begin{itemize}\small
        \item[\textbf{P1}] ¿Tienen las interacciones digitales poder predictivo real sobre el comportamiento electoral en Colombia?
        \item[\textbf{P2}] ¿Es posible reconstruir matemáticamente el historial de interacciones sin haber recolectado datos en tiempo real?
        \item[\textbf{P3}] ¿Qué métrica digital es el predictor más informativo y robusto?
        \item[\textbf{P4}] ¿Cuántas contiendas de entrenamiento son necesarias para que el modelo alcance desempeño estable?
      \end{itemize}
    \end{column}
    \begin{column}{0.48\textwidth}
      {\small\textbf{\color{barBlue}Preguntas metodológicas (PM)}}
      \begin{itemize}\small
        \item[\textbf{PM1}] ¿Es posible replicar el SOMEN-DC en una elección no tratada por sus autores?
        \item[\textbf{PM2}] ¿Se confirman empíricamente las tres limitaciones estructurales del SOMEN-DC?
        \item[\textbf{PM3}] ¿Es posible reconstruir el historial de interacciones con precisión suficiente para construir un dataset de entrenamiento retroactivo?
        \item[\textbf{PM4}] ¿Es el modelo transferible a un formato electoral diferente al de entrenamiento?
      \end{itemize}
    \end{column}
  \end{columns}
\end{frame}

% DIAPOSITIVA 6: Diseño general
\begin{frame}{Diseño empírico: tres etapas}
  \begin{columns}[T]
    \begin{column}{0.42\textwidth}
      \vspace{0.3em}
      \begin{center}
      \begin{tikzpicture}[node distance=0.5cm]
        \tikzstyle{box}=[rectangle, rounded corners=3pt, draw=barBlue,
                         fill=tableBand, minimum width=3.8cm, minimum height=0.7cm,
                         text width=3.6cm, align=center, font=\small\bfseries]
        \tikzstyle{lbl}=[font=\scriptsize, text=subtleGray, align=center,
                         text width=3.6cm]
        \tikzstyle{arr}=[->, thick, color=barBlue]

        \node[box] (bol) {Etapa 1: Bolivia};
        \node[lbl, below=0.15cm of bol] (lbol) {Replicación SOMEN-DC\\PM1, PM2};
        \node[box, below=0.4cm of lbol] (cr) {Etapa 2: Costa Rica};
        \node[lbl, below=0.15cm of cr] (lcr) {Restaurador Weibull\\P2, PM3};
        \node[box, below=0.4cm of lcr] (col) {Etapa 3: Colombia};
        \node[lbl, below=0.15cm of col] (lcol) {Modelo ELA-NOM\\P1, P3, P4};
        \node[box, below=0.4cm of lcol] (val) {Validación 2026};
        \node[lbl, below=0.15cm of val] (lval) {Transferibilidad\\PM4};

        \draw[arr] (bol) -- (cr);
        \draw[arr] (cr) -- (col);
        \draw[arr] (col) -- (val);
      \end{tikzpicture}
      \end{center}
    \end{column}
    \begin{column}{0.54\textwidth}
      \textbf{Variable objetivo:}
      \[
        \hat{y}_c = f\!\left(\mathbf{x}_c\right) \in (0,1)
        \quad\text{con}\quad
        \sum_{j \in \mathcal{C}_e} \hat{y}_{c_j} = 1
      \]
      \small
      $f\!: \mathbb{R}^d \to (0,1)$, escalar por candidato.
      La restricción de suma se aplica por normalización posterior.

      \vspace{0.5em}
      \textbf{Supuesto central:}
      \begin{itemize}
        \item Los parámetros de $f$ se aprenden con \textbf{múltiples} elecciones pasadas $\mathcal{E}$
        \item Se aplican \textbf{sin reentrenamiento} a elecciones nuevas
        \item Variable objetivo: \kw{cuota electoral observada} (no intención de voto latente)
      \end{itemize}
    \end{column}
  \end{columns}
\end{frame}
